\section{Darstellung von Zahlen}
\label{sec:fundamentals_number_representation}

Im Alltag ist die \emph{Dezimaldarstellung} die Standardschreibweise für
Zahlen. In dieser Darstellung besitzt jede Position einer Ziffer einen festen
Stellenwert: Tausender, Hunderter, Zehner und Einer sowie nach dem Komma Zehntel,
Hundertstel, Tausendstel, usw...

{\centering
  \begin{tikzpicture}[x=2.1cm,y=1cm,font=\small]
  \foreach \x/\digit/\power/\name in {
    0/1/{10^3}/{\Lang{Tausender}{thousands}},
    1/1/{10^2}/{\Lang{Hunderter}{hundreds}},
    2/5/{10^1}/{\Lang{Zehner}{tens}},
    3/3/{10^0}/{\Lang{Einer}{ones}},
    4/5/{10^{-1}}/{\Lang{Zehntel}{tenths}},
    5/3/{10^{-2}}/{\Lang{Hundertstel}{hundredths}},
    6/7/{10^{-3}}/{\Lang{Tausendstel}{thousandths}}} {
    \node[font=\footnotesize,align=center,text width=1.9cm] at (\x,1.65)
      {\name};
    \node[draw=black!35,fill=black!2,minimum width=1.8cm,
      minimum height=0.9cm,font=\Large\bfseries] at (\x,0.85) {$\digit$};
    \node[text=black!60] at (\x,0.1) {$\power$};
  }

  \node[font=\footnotesize,text=THMBlue] at (3.5,2.15)
    {\Lang{Komma}{decimal point}};
  \draw[->,THMBlue,line width=0.6pt] (3.5,1.95) -- (3.5,1.18);
  \node[font=\Large\bfseries,text=THMBlue] at (3.5,0.85) {$,$};

  \draw[THMBlue,line width=0.8pt] (-0.44,-0.35) -- (3.44,-0.35);
  \draw[THMBlue,line width=0.8pt] (3.56,-0.35) -- (6.44,-0.35);
  \node[font=\scriptsize,text=THMBlue] at (1.5,-0.65)
    {\Lang{Vorkommastellen}{integer part}};
  \node[font=\scriptsize,text=THMBlue] at (5,-0.65)
    {\Lang{Nachkommastellen}{fractional part}};
\end{tikzpicture}

  \par}

Bei sehr großen und sehr kleinen Zahlen ist diese Darstellung jedoch umständlich
und unübersichtlich, da um die Größenordnung einzuschätzen alle Stellen abgezählt
werden müssen. Abhilfe schafft die dezimale Gleitkommadarstellung.

\begin{DefinitionBox}{Dezimale Gleitkommadarstellung}{fundamentals_decimal_floating_point}
  \noindent
  \begin{minipage}[c]{0.42\linewidth}
    Jede Zahl $x\neq0$ lässt sich normalisiert in Gleitkommadarstellung
    schreiben.

    \smallskip
    Es gilt $1\leq |m|<10$ und $e\in\Z$. Die Mantisse beginnt daher mit genau
    einer von null verschiedenen Ziffer.
  \end{minipage}\hfill
  \begin{minipage}[c]{0.54\linewidth}
    \centering
    \begin{tikzpicture}[
    font=\normalsize,
    formula/.style={font=\Large},
    callout/.style={font=\normalsize,align=center},
    pointer/.style={-{Latex[length=2.2mm]},THMBlue,line width=0.65pt}
  ]
  \node[formula] (lhs) at (-1.65,0) {$x={}$};
  \node[formula] (mantissa) at (-0.65,0) {$m$};
  \node[formula] (times) at (0,0) {$\cdot$};
  \node[formula] (base) at (0.7,0) {$10$};
  \node[font=\normalsize,anchor=south west]
    (exponent) at ([xshift=-0.05cm,yshift=-0.05cm]base.north east) {$e$};

  \node[callout] (mantissa-label) at (-0.65,-0.8)
    {\Lang{Mantisse}{mantissa} $m$};
  \node[callout] (exponent-label) at (1.75,0.75)
    {\Lang{Exponent}{exponent} $e$};

  \draw[pointer] (mantissa-label.north) -- (mantissa.south);
  \draw[pointer] (exponent-label.south west) -- (exponent.north east);
\end{tikzpicture}

  \end{minipage}
\end{DefinitionBox}

\begin{ExampleBox}{Darstellungen im Vergleich}{fundamentals_floating_point_comparison}
  \renewcommand{\arraystretch}{1.15}
  \begin{tabularx}{\linewidth}{@{}>{\raggedright\arraybackslash\bfseries}p{3.9cm}
      >{\raggedright\arraybackslash}X
      >{\raggedright\arraybackslash}p{4.4cm}@{}}
    \toprule
    Größe & \textbf{Dezimaldarstellung}
          & \textbf{\shortstack[l]{Gleitkomma-\\darstellung}}\\
    \midrule
    {}
      & $1153{,}537$
      & $1{,}153537\cdot10^3$\\[0.35em]
    Masse der Erde
      & $5\,972\,200\,000\,000\,000\,000\,000\,000\,\mathrm{kg}$
      & $5{,}9722\cdot10^{24}\,\mathrm{kg}$\\[0.35em]
    Leiterbahnbreite
      & $0{,}000\,000\,003\,\mathrm m$
      & $3\cdot10^{-9}\,\mathrm m$\\
    \bottomrule
  \end{tabularx}
  \renewcommand{\arraystretch}{1}
\end{ExampleBox}

\begin{KeyPointBox}[breakable=false]{Vorteile der Gleitkommadarstellung}{fundamentals_floating_point_advantages}
  \begin{enumerate}[label=\alph*),leftmargin=2em,itemsep=0.2em]
    \item Sehr große und sehr kleine Zahlen werden kompakt dargestellt.
    \item Der Exponent zeigt die Größenordnung und erleichtert Vergleiche.
    \item Potenzgesetze vereinfachen Multiplikationen unterschiedlicher
      Größenordnungen, z.B.
      \[
        (2\cdot10^3)\cdot(4\cdot10^{-5})=(2\cdot4)\cdot10^{3-5}=8\cdot10^{-2}.
      \]
    \item Durch Runden der Mantisse auf die benötigte Genauigkeit entfallen
      nicht-aussagekräftige Ziffern, z.B.
      \[
        6\,724\,891\,356\,000\,000\,000\,000\,000 \approx 6{,}72\cdot10^{24}.
      \]
  \end{enumerate}
\end{KeyPointBox}
Aus diesem Gründen ist die Gleitkommadarstellung in Naturwissenschaft und Technik sehr verbreitet.
